\documentclass[11pt,letterpaper]{article}
\usepackage[margin=1in]{geometry}
\usepackage[T1]{fontenc}
\usepackage{amsmath,amsthm,booktabs,array,longtable,graphicx}
\usepackage{newtxtext,newtxmath}
\usepackage[round,authoryear]{natbib}
\usepackage{setspace,xr-hyper,xurl}
\usepackage[hidelinks]{hyperref}
\usepackage{fancyhdr}
\onehalfspacing\setlength{\emergencystretch}{2em}
\newtheorem{theorem}{Theorem}[section]
\newtheorem{proposition}[theorem]{Proposition}
\newtheorem{lemma}[theorem]{Lemma}
\newtheorem{corollary}[theorem]{Corollary}
\theoremstyle{definition}\newtheorem{remark}[theorem]{Remark}
\newcommand{\R}{\mathbb R}\newcommand{\E}{\mathbb E}
\newcommand{\argmax}{\operatorname*{arg\,max}}
\newcommand{\argmin}{\operatorname*{arg\,min}}
\pagestyle{fancy}\fancyhf{}\fancyhead[L]{\small Finite-Catalog Resource Allocation}
\fancyhead[R]{\small Anonymous}\fancyfoot[C]{\thepage}\setlength{\headheight}{14pt}

% BEGIN retained/current source: revisions/or-r52-resource-path-20260925/generated/metrics.tex
\newcommand{\FineMaxGap}{0.000500}
\newcommand{\FineMinTime}{3.523}
\newcommand{\FineMaxTime}{19.762}
\newcommand{\FineMaxRelative}{1.63\%}
\newcommand{\FineMinPacked}{0.216}
\newcommand{\FineMaxPacked}{6.628}
\newcommand{\FineMaxDenBits}{553}
\newcommand{\AllResourceExact}{31}
\newcommand{\GroupNOneComplete}{2}
\newcommand{\MIPMaxSeconds}{0.038}
\newcommand{\MIPMaxGap}{$5.55\,10^{-17}$}

% END source


% BEGIN retained/current source: revisions/or-r54-exact-price-path-20260926/generated/metrics.tex
\newcommand{\PriceExact}{10}
\newcommand{\PriceTolerance}{19}
\newcommand{\PriceLimited}{1}
\newcommand{\PriceChecked}{30}
\newcommand{\PriceMinBytes}{541}
\newcommand{\PriceMaxBytes}{5895}
\newcommand{\AnytimeExact}{10}
\newcommand{\ForcedAnchors}{45}
\newcommand{\HardnessSubsets}{1692}
\newcommand{\HardnessBooks}{3004}

% END source

\hypersetup{pdftitle={Finite-Catalog Resource Allocation: A Tariff-Sensitive Complexity Frontier},pdfauthor={Anonymous}}
\externaldocument{.build/r63/electronic_companion_refs}[electronic_companion.pdf]
\begin{document}\hypersetup{pageanchor=false}
\begin{titlepage}\centering{\Large\bfseries Finite-Catalog Resource Allocation: A Tariff-Sensitive Complexity Frontier\par}
\vspace{0.25in}Anonymous manuscript for Operations Research --- R63\par\vspace{0.2in}
\begin{minipage}{\textwidth}\textbf{Abstract.} A service provider must allocate an accepted aggregate promise while installing a small shared catalog of terminal commands. Individual expected caps and realization ceilings constrain implementation, and every installed command incurs a fixed charge. We derive an exact resource-path representation and establish a tariff-sensitive complexity frontier. With common linear rewards and free preliminary service, arbitrary command-specific charges remain parameterized-hard in the command allowance and cap count. A capacity dynamic program instead gives exact fixed-parameter tractability in the number of exceptions to a standard fee, and uniform fees permit polynomial-time optimization without a resource grid. We extend this result to near-standard tariffs through an explicit original-policy perturbation certificate and a minimum-distortion sparse-exception projection. Exact fee frontiers identify changes in installed commands and service allocation. Conservation, additive approximation, and same-book price support provide complementary guarantees for broader regimes. Controlled scaling and adversarial studies retain every interrupted request, separate numerical bounds from rational certification, and measure optimization, proof size, memory, and independent checking. An additional exhaustive diagnostic distinguishes genuine root support from positive price gaps.
\par\vspace{0.15in}\textbf{Keywords:} resource allocation; dynamic programming; parameterized complexity.
\par\vspace{0.1in}\textbf{Subject classifications:} Programming: resource allocation; Dynamic programming: finite-catalog design.
\par\vspace{0.1in}\textbf{Area of review:} Optimization.\end{minipage}\vfill September 27, 2026\end{titlepage}
\hypersetup{pageanchor=true}
\section{Introduction}
A finite command catalog creates a distinction between the resources allocated to a decision and the number of commands used to implement it. A provider may deliver different expected obligations to different customer histories while installing only a small common set of terminal commands. We call the installed set a command book. Each installed command incurs a fixed charge. Individual expected caps and ceilings on every realized obligation restrict how an accepted aggregate promise can be implemented. Choosing expected targets first and a catalog afterward can therefore miss both a discrete cost and a feasibility constraint.

This paper establishes a tariff-sensitive complexity frontier for joint resource allocation and implementation. The main contrast holds even when rewards are linear and preliminary service is free. With arbitrary command-specific opening fees, exact optimization remains parameterized-hard in both the command allowance and the number of distinct expected caps. With a standard fee and only a few exceptional commands, however, we obtain an exact fixed-parameter algorithm whose polynomial exponent does not depend on the number of exceptions. Uniform fees give a polynomial-time algorithm with no discretization of capacities or the accepted promise. Thus a small command allowance does not by itself make the problem tractable, whereas a simple installation tariff can do so without restricting the physical heterogeneity.

The positive result follows from a capacity representation inside the original model. A selected book's gross value depends only on its attainable aggregate terminal capacity. This capacity telescopes along consecutive selected commands. Conditional on the installed exceptional commands and the book size, the total opening charge is fixed; it is then sufficient to retain one maximum-capacity label at each state. A recovered label is not merely a numerical value: it supplies terminal lotteries and preliminary service that meet the original promise and every individual constraint. The same tables compute the provider's exact decision frontier as a uniform installation fee changes. The dynamic-programming principle is familiar; its applicability here follows from the model-specific capacity and fixed-charge identities.

The exact selected-boundary representation also supports the general concave model. We characterize capacity conservation on acyclic graphs and extend the recognition criterion to paths with an explicit edge allowance. Complementing resources merges feasible starting points into a single deficit state. Rounding and repair preserve the selected path, its command count and the accepted promise. These results retain the original-feasible additive guarantee even where the tariff algorithm is inapplicable. A corridor-rehabilitation model illustrates the conservation and repair mechanism outside renewal commands, with its physical quantities and separability assumptions stated explicitly in the companion.

Price-based certificates and resource approximation play different roles. We characterize exactly when one price-active book supports the accepted promise, and bound the residual gap by the distance to that book's maximizing target interval. This explains why price methods can be strong on root-supported instances without implying that their general search trees are small. A uniform positive installation fee can still create a price gap even though the tariff algorithm solves the problem exactly. The earlier deficit implementation's unfavorable timing results remain part of the evidence; an accuracy guarantee does not establish practical speed superiority.

The computational study preserves the earlier evidence and adds an expanded frozen scaling and adversarial design. It varies the number of fee exceptions separately from physical size, records memory and proof costs, and retains every interrupted request. A separate exhaustive diagnostic identifies when the whole price-active family supports the promise. Near-standard fees are not treated as exact uniform fees: we bound the value lost by simplifying the tariff and evaluate the returned policy under its original charges.

\subsection{Relationship to prior work}
Lagrangian constrained-path bounds, label search and resource-state approximation are established tools \citep{HandlerZang1980,BeasleyChristofides1989,Fisher1981,Hassin1992,LorenzRaz2001,PuglieseGuerriero2013}. Contemporary bucket-graph labeling organizes dominance and reduced-cost elimination for routing subproblems \citep{Sadykov2021}; heuristic-guided constrained search addresses large networks and multiple resource limits \citep{Ahmadi2025}. Those methods are not displaced by our deficit dynamic program. Our graph edges instead carry allocatable continuous resources, and our claims concern an exact representation of individual obligations, a recognizable conservation restriction, and recovery of an original equality. Generic path algorithms do not establish those model identities. No timing comparison against a routing implementation is inferred from our allocation experiments.

Continuous allocation on a fixed book belongs to separable resource allocation \citep{Patriksson2008}. Fixed-charge nonlinear allocation also admits perspective reformulations and branch-and-Benders-cut approaches \citep{Yang2025}. These are important formulation and algorithmic antecedents, not results that opening fees generally become easy. Our positive theorem exploits the additional ordered-capacity identity and the fact that conditioning on exceptional commands makes a given book size have a common charge. Its hardness counterpart concerns parameterized complexity, rather than the presence of binary variables alone. Parameterized knapsack classifications distinguish structural parameters, numerical values and their encodings \citep{Gurski2019}; their bounds do not transfer to our shared-book constraints without a reduction. We supply that reduction directly inside the original policy model.

Small distributional support is not a small shared installation catalog. Moment-set extreme-point results bound individual supports \citep{Winkler1988}; they do not charge for their union. Auction menu complexity additionally imposes incentive constraints absent here \citep{HartNisan2013,Babaioff2016}. Choice-based assortment optimization changes demand through the offered set \citep{Rusmevichientong2010}, whereas our history probabilities are fixed and terminal probabilities are controlled. Our reward perturbation certificate holds feasibility inputs fixed; it is not a statistical distributional-robustness claim \citep{EsfahaniKuhn2018}.

Finally, concave convolution uses classical matrix-searching structure \citep{Aggarwal1987}; the finite-support proof permits holes created by merged sources. General medoid selection and one-dimensional weighted-median clustering are antecedent algorithms, not independent novelty claims \citep{Arya2004,Wu1991}. Their role is to make the model's oscillation-based grouping bound operational. Table~\ref{tab:frontier60} separates these antecedents from the results established here. The paper's scientific contribution is the exact representation and tariff-sensitive frontier, with conservation, approximation and certification supplying implementable consequences.
\section{Finite-Catalog Renewal Design}\label{sec:model52}
There are $k$ histories. History $j$ has probability $\pi_j>0$, expected-obligation cap $b_j\in[0,1]$, and realization ceiling $\tau_j\geq b_j$, with $\sum_j\pi_j=1$. Write $\bar b=\sum_j\pi_jb_j$. Intermediate service has cost $h_j:[0,b_j]\to\mathbb R_+$, continuous, convex and nondecreasing, with $h_j(0)=0$. The initial common terminal reward $f:[0,1]\to\mathbb R$ is strictly increasing and concave. The catalog is $\mathcal A=\{a_1<\cdots<a_N\}\subset[0,1]$. Opening command $a_i$ costs $\rho_i\geq0$.

A nonempty selected book $c=(c_1<\cdots<c_s)$ has $s\leq m$. After observing $j$, the provider selects deterministic pre-draw service $y_j\geq0$ and a lottery $C_j$ supported on $c$. The problem is
\begin{align}
 V^*(B)&=\max_{c,y,C}\left\{\sum_j\pi_j\bigl[\mathbb E f(C_j)-h_j(y_j)\bigr]-\sum_{z\in c}\rho(z)\right\},\label{eq:original52}\\
 t_j&=y_j+\mathbb E C_j\leq b_j,\qquad y_j+C_j\leq\tau_j\quad\text{almost surely},\qquad \sum_j\pi_jt_j=B.\label{eq:constraints52}
\end{align}
Here $B\in[0,\bar b]$ is accepted, rather than optimized away. The common-risk specialization is $\tau_j=b_j+\delta$. Expected and realized obligations are different restrictions. The terminal symbol does not carry a free history index, inherited service tier, or shared random seed. The cardinality $m$ measures the command alphabet, not total implementation storage: targets, encoder logic, probability precision, program constants, and certificates are additional resources. The mathematical results below concern the stated finite catalog; continuous-design results have a separate approximation error.

For a feasible book, let $d_j=\max(c\cap[0,\tau_j])$, and let $F_{j,c}$ interpolate $f$ linearly on this eligible prefix. A singleton prefix has the constant value $f(d_j)$ at its only point.
\begin{proposition}[Canonical response]\label{prop:canonical52}
A book is feasible if and only if $c_1\leq\min_jb_j$ and $c_1\leq B\leq\bar b$. At target $c_1\leq t\leq b_j$, its optimal response is
\begin{equation}\label{eq:canonical52}
 v_{j,c}(t)=F_{j,c}(\min\{t,d_j\})-h_j((t-d_j)_+).
\end{equation}
It is attained by $y_j=(t-d_j)_+$ and an adjacent eligible lottery with mean $\min(t,d_j)$. The response is concave in $t$.
\end{proposition}
\begin{proof}
Every feasible command satisfies $C_j\leq\tau_j$ since $y_j\geq0$. At mean $\mu=\mathbb E C_j$, concavity on the ordered support bounds the reward by $F_{j,c}(\mu)$. The payoff is at most $F_{j,c}(\mu)-h_j(t-\mu)$ for $c_1\leq\mu\leq\min(t,d_j)$. Both terms are nondecreasing in $\mu$, so the upper endpoint maximizes this bound. When $t\leq d_j$, adjacent interpolation has zero intermediate service and support below $\tau_j$. When $t>d_j$, use command $d_j$ surely and service $t-d_j$; the realized sum is $t\leq b_j\leq\tau_j$. Thus the bound is attained without weakening either restriction. Interpolation slopes decrease and are nonnegative; the tail slopes of $-h_j$ are nonpositive and nonincreasing. This proves concavity. Feasible targets fill the entire weighted interval $[c_1,\bar b]$, proving necessity and sufficiency.
\end{proof}

The algorithmic specialization is
\begin{equation}\label{eq:quadratic52}
 f(x)=rx-\tfrac12 qx^2,\quad r\geq q\geq0,\quad r>0,\qquad h_j(y)=\tfrac12\gamma_jy^2,\quad \gamma_j\geq0.
\end{equation}
All probabilities, caps, ceilings, catalog points, shape parameters, opening charges, and requested accuracies are rationally encoded in exact-arithmetic claims. Optimization identities allow real nonnegative charges; bit-complexity assertions do not assign a finite encoding to arbitrary real numbers. A common Lipschitz constant for the canonical responses is
\begin{equation}\label{eq:L52}L=\max\{r,\max_j\gamma_j\}>0.\end{equation}
Continuity of the cost on its compact domain makes the stated extrema attained; every rational quadratic instance already satisfies this regularity. A fixed-book allocation can be solved by filling positive chord slopes in decreasing order, then zero-cost service capacity, then quadratic service tails. The service breakpoints are $\gamma_j(b_j-d_j)_+$. This takes $O(ks+k\log(k+1))$ rational operations with sorted marginal events. Alternatively, a direct supporting-price sweep verifies
\begin{equation}\label{eq:fixeddual52}
 \sum_j\pi_jv_{j,c}(t_j)=\lambda B+\sum_j\pi_j\max_{c_1\leq t\leq b_j}\{v_{j,c}(t)-\lambda t\}.
\end{equation}
Eligible knots, caps, and clipped stationary tail points suffice to check each scalar maximum. This certificate proves optimality for that book, not over all books.
\subsection{Infeasible inputs and certificate scope}\label{sec:infeasible59}
For well-formed primitives and a nonempty catalog, feasibility is decided before an incumbent is initialized. If $m<1$, $B<0$, $B>\bar b$, or $a_1>\min\{B,\min_jb_j\}$, no feasible book exists. Necessity follows from nonemptiness, the aggregate caps and the minimum possible command; sufficiency follows by using the singleton $a_1$ and filling deterministic service to the accepted promise. An empty catalog is likewise infeasible. A malformed probability vector, unsorted catalog or invalid primitive is instead an input error and must not be reported as an optimization certificate.

An infeasibility record identifies the failed inequality and carries the external input digest. It has no incumbent, lower policy or finite objective interval. An independent checker recomputes the probability-weighted cap, minimum cap and minimum catalog point from the supplied original input. Validating a scalar inequality is sufficient; search is unnecessary. The current solver modules retain their legacy exception on infeasible requests, whereas \texttt{feasibility59.py} provides this explicit preflight record and checker. The new panel is generated feasible and does not count preflight rejection as an optimization success.
\paragraph{Units and loss.} Obligations and caps use one normalized resource unit. Rewards, service costs and installation charges share a net-value unit. An additive tolerance is stated in that value unit; rescaling every payoff and the tolerance by the same positive factor preserves its decision meaning. The present theory makes no empirical currency calibration.
\subsection{An operational interpretation of the tariff restriction}\label{sec:institution63}
One interpretation is a provider that has already accepted a service-credit obligation. Histories index observable eligibility classes; preliminary service and a randomized terminal credit jointly discharge the promised obligation. The expected cap limits the average obligation to a class, while the realization ceiling limits any delivered outcome. A command book is the set of terminal credit levels configured before individual assignments. A common opening charge represents a standardized setup procedure; exceptional levels require additional integration or qualification. This is a model interpretation, not a claim that a particular provider has supplied or adopted these inputs.

The accepted aggregate promise is a commitment to be met exactly. It is not obtained by changing an upper spending budget into an equality. The operational decision is therefore how to implement an already contracted obligation, including its fixed installation charges. The unrestricted model permits concave terminal benefits and costly preliminary service. The tariff theorem identifies a specific physical regime---a common linear benefit and zero incremental preliminary-service cost---in which fee structure alone changes computational tractability. Heterogeneous caps, ceilings, probabilities, catalog spacing, and an unrestricted command allowance remain present. Near-standard tariffs are treated by the explicit perturbation certificate in Section~\ref{sec:robust63}, rather than by silently rounding measured charges.
\section{Constructive Selected-Boundary Decomposition}\label{sec:boundaries52}
This section assumes the common reward $f$ in Section~\ref{sec:model52}. For a selected book $c=(u_1<\cdots<u_s)$, write $a=u_1$ and $d_j=\max\{u_i:u_i\leq\tau_j\}$.
\begin{lemma}[Global terminal-first allocation]\label{lem:terminal52}
An optimal fixed-book canonical allocation cannot have both $t_i>d_i$ and $t_j<\min(b_j,d_j)$.
\end{lemma}
\begin{proof}
Choose $0<\delta\leq\min\{t_i-d_i,\pi_j[\min(b_j,d_j)-t_j]/\pi_i\}$. Reduce $t_i$ by $\delta$ and increase $t_j$ by $\pi_i\delta/\pi_j$. The weighted change is zero. The donor's terminal reward is unchanged and its service cost weakly decreases; the recipient's strictly increasing terminal interpolant improves. Both responses remain feasible by Proposition~\ref{prop:canonical52}, contradicting optimality.
\end{proof}
For selected neighbors $u<v$, define
\begin{align}
 \delta_j(u,v)&=\mathbf1_{\{\tau_j\geq v\}}[\min(b_j,v)-u]_+,
 &T_{uv}&=\sum_j\pi_j\delta_j(u,v),\label{eq:terminalarc52}\\
 J_{uv}&=\{j:u\leq\tau_j<v\},
 &Q_{uv}&=\sum_{j\in J_{uv}}\pi_j(b_j-u)_+.\label{eq:exit52}
\end{align}
The capped service value on an exit set is
\begin{equation}\label{eq:servicearc52}
 H_{uv}(y)=\min\left\{\sum_{j\in J_{uv}}\pi_jh_j(z_j):
 0\leq z_j\leq(b_j-u)_+,\ \sum_{j\in J_{uv}}\pi_jz_j=y\right\},\quad 0\leq y\leq Q_{uv}.
\end{equation}
Empty sets have zero capacity and zero value at zero. Let $\sigma_{uv}=[f(v)-f(u)]/(v-u)$ and
\begin{equation}\label{eq:arc52}
 \Phi_{uv}(x)=\sigma_{uv}\min(x,T_{uv})-H_{uv}((x-T_{uv})_+),
 \qquad 0\leq x\leq T_{uv}+Q_{uv}.
\end{equation}
For the terminal edge $(u,\dagger)$, use $T_{u\dagger}=0$, $J_{u\dagger}=\{j:\tau_j\geq u\}$, and $\Phi_{u\dagger}=-H_{u\dagger}$. The functions $\Phi$ are concave. Under the quadratic primitives, they are $L$-Lipschitz with $L$ as in \eqref{eq:L52}.

\begin{lemma}[Nested terminal layers]\label{lem:nested54}
Start each history at target $a$. Fill the selected terminal layers in increasing edge order. Once every preceding layer is full, each history with $\delta_j(u,v)>0$ is at target $u$. Every weighted amount $x\in[0,T_{uv}]$ in the next layer has a feasible individual allocation. After all terminal layers are full, history $j$ is at $\min(b_j,d_j)$; its remaining service capacity belongs to exactly one exit set.
\end{lemma}
\begin{proof}
A history with positive increment on $(u,v)$ has $b_j>u$ and $\tau_j\geq v$. It is eligible for all earlier selected commands, and the increments of those earlier layers telescope to $u-a$. This proves the first assertion, including a history whose cap will truncate the current layer.

For any desired $x$, initialize residual $r=x$ and visit histories in any fixed order. Give history $j$ increment $e_j=\min\{\delta_j(u,v),r/\pi_j\}$ and reduce $r$ by $\pi_je_j$. Since $x\leq\sum_j\pi_j\delta_j(u,v)$, the final residual is zero. Each affected target lies in $[u,\min(b_j,v)]$ and is implemented by the adjacent lottery on $u,v$ with upper probability $e_j/(v-u)$. It has no intermediate service and both realizations are below $\tau_j$. Histories with zero increments keep their previously feasible responses.

If a layer is partial, stop: its predecessors are full and its successors unused. If every layer is full, the increments for history $j$ telescope to $\min(b_j,d_j)-a$. It exits exactly after $d_j$, at an ordinary crossing edge or at the terminal edge. Thus its remaining capacity is $(b_j-d_j)_+$ and appears in exactly one $H$. Any feasible service vector on the exit sets implements target $\min(b_j,d_j)+z_j$. A positive $z_j$ requires $b_j>d_j$, and its implementation is the sure command $d_j$ plus service $z_j$; their sum is at most $b_j\leq\tau_j$. Zero-capacity histories require no service. The exit sets are disjoint, so these implementations combine without conflicting allocations.
\end{proof}

\begin{lemma}[Implementable rearrangement]\label{lem:rearrange54}
For a fixed selected path, any feasible arc-resource vector can be rearranged, with unchanged total resource and no smaller payoff, into full earlier terminal layers, at most one partial terminal layer, and service only after all terminal layers are full. That rearranged payoff is attained by original history-level policies.
\end{lemma}
\begin{proof}
Split the resource on each arc into its terminal and service parts in \eqref{eq:arc52}. If any terminal capacity remains unused while some service is positive, transfer the smaller of the available terminal capacity and a positive service amount to that terminal layer. The terminal gain is positive, while reducing service weakly lowers cost. Repeat until there is no such pair. This operation never changes charges or the sum of resources.

Next, if an earlier terminal layer is unfinished and a later one is positive, transfer the smaller of the unfinished capacity and the later mass. Selected chord slopes are positive and nonincreasing, so payoff cannot decrease. Filling layers successively yields the asserted form after finitely many exhausted donors or filled recipients. Each arc still respects its capacity: terminal mass is at most $T_e$, and service, when present, is at most $Q_e$ after that arc's terminal part is full. Allocate any remaining service by minimizers of the corresponding $H_e$; these minimizers exist by compactness and continuity. Lemma~\ref{lem:nested54} implements the resulting terminal pattern and the disjoint service allocations. Its gross payoff is precisely $f(a)+\sum_e\Phi_e(x_e)$ for the rearranged vector. An arbitrary original vector need not be implemented at its unrearranged payoff; the weak improvement is the comparison needed for exact optimal values.
\end{proof}

\begin{theorem}[Exact selected-boundary resource path]\label{thm:path52}
Let $P(c)$ contain the successive selected edges and the terminal edge of a feasible book. Then
\begin{align}
 \sum_{e\in P(c)}(T_e+Q_e)&=\bar b-a,\label{eq:capacityidentity52}\\
 V_c(B)&=f(a)-\sum_{z\in c}\rho(z)+
 \max_{\substack{0\leq x_e\leq T_e+Q_e\\\sum_e x_e=B-a}}\sum_{e\in P(c)}\Phi_e(x_e).\label{eq:pathidentity52}
\end{align}
Maximizing over paths with at most $m$ selected commands solves the original joint problem exactly. Reconstruction fills terminal layers in order, uses the greedy partial-layer allocation of Lemma~\ref{lem:nested54}, and then solves the disjoint capped service allocations. The resulting policy obeys the original promise, expected caps, realization ceilings, and command budget.
\end{theorem}
\begin{proof}
For history $j$, the terminal increments sum to $\min(b_j,d_j)-a$ and its unique service capacity is $(b_j-d_j)_+$. Their sum is $b_j-a$, proving \eqref{eq:capacityidentity52} after weighting. In an original optimum, Lemma~\ref{lem:terminal52} places service after terminal capacity. Among terminal layers, exchanging equal weighted amounts orders the common decreasing chord slopes; Lemma~\ref{lem:nested54} justifies the individual feasibility of that ordering. If all terminal capacity is full, the exit-set cost minimizations describe all remaining feasible service allocations. Hence an original optimum supplies an arc allocation of the same value. Conversely, Lemma~\ref{lem:rearrange54} implements a weakly improved version of every feasible arc allocation. The two comparisons yield equality of the attained maxima. Each selected charge is subtracted once, and the terminal edge does not open an additional command.
\end{proof}

\begin{corollary}[Endogenous eligibility]\label{cor:endogenous52}
A book of $s$ commands has at most $s$ nonempty selected-prefix classes. Inserting unselected candidates changes neither its arc functions nor its fixed-book optimum.
\end{corollary}
\begin{proof}
The nonempty exit sets partition histories according to their last eligible selected command. Their definitions and terminal increments depend on the selected endpoints, not intervening unused candidates.
\end{proof}
\section{The Limited-Menu Complexity Frontier}\label{sec:param59}
Let $q_b=|\{b_j:j=1,\ldots,k\}|$ denote the number of distinct expected caps; the subscript distinguishes this count from a quadratic curvature. For a fixed command budget, enumerating all feasible books is polynomial with an exponent depending on that budget. Under a common reward, Theorem~\ref{thm:capcount56} reduces the useful book size to $\min\{m,2q_b\}$, and its coefficient two is tight. The following result identifies why neither parameter generally removes the exponent dependence.

\subsection{Encoded decision problem and the meaning of the parameters}\label{sec:decision60}
For the complexity results, an instance is an ordered list of rational catalog levels and nonnegative opening charges, positive rational probabilities summing to one, rational expected caps and realization ceilings, a rational accepted promise and value threshold $z$, an integer command allowance $m\geq1$, and the rational quadratic reward and service coefficients of Section~\ref{sec:model52}. Each rational is represented by a signed binary numerator and a positive binary denominator in lowest terms. Sorting and validation are polynomial-time preprocessing; duplicate catalog entries are consolidated at their least opening charge. The decision language asks whether an original feasible policy using a nonempty book of at most $m$ commands has net value at least $z$. Invalid or infeasible inputs are no-instances. Denote the total encoding length by $L$ and the number of distinct reduced rational caps by $q_b$.

Theorem~\ref{thm:wone59} is a parameterized many-one lower bound for this language with parameter $m+q_b$, already in the common linear-reward, zero-service-cost subclass. Hence it also excludes an algorithm $f(m)L^{O(1)}$ or $f(q_b)L^{O(1)}$, unless FPT equals W[1]: either such algorithm would be fixed-parameter in the sum. It does not assert W[1] membership or completeness, an ETH lower bound on the exact exponent, strong NP-hardness, or a kernel lower bound. The fixed-cap reconstruction and exact-lattice algorithms address different restrictions. In particular, numerical denominators cannot be replaced by their bit lengths in a pseudo-polynomial complexity expression. The positive result in Section~\ref{sec:tariff60} instead removes denominator dependence from the operation count by controlling opening-fee exceptions.
\begin{theorem}[Parameterized hardness in menu size and cap count]\label{thm:wone59}
The rational finite-catalog decision problem is $\mathrm{W}[1]$-hard parameterized by $m+q_b$. Hardness holds with the fixed common reward $f(x)=x$, zero service costs, uniform strictly positive probabilities, $\tau_j=b_j$, nonnegative rational opening charges, and $m=q_b$. Consequently, unless $\mathrm{FPT}=\mathrm{W}[1]$, no exact algorithm has running time $F(m,q_b)\mathcal L^c$ for an arbitrary computable $F$ and a constant $c$ independent of the parameters, where $\mathcal L$ is binary input length. The same conclusion holds for parameterization by either $m$ or $q_b$ alone.
\end{theorem}

The proof has two explicit stages. A group-sum encoding carries the consistency requirements of a multicolored clique. Dedicated baseline histories then make those group choices mandatory in the original renewal problem. Unlike a free augmentation under a nonbinding budget, every baseline must actually be installed in a threshold-achieving book.

\begin{lemma}[Group-sum encoding]\label{lem:groupsum59}
From an instance of multicolored clique with $r$ color classes, one can construct $g=r+\binom r2$ nonempty groups of positive integers $A_1,\ldots,A_g$ and a positive integer $W$, of total binary length polynomial in the graph description, such that a selection of at most one integer from each group sums to $W$ if and only if the graph has a multicolored clique. Every such sum uses exactly one integer from every group.
\end{lemma}
\begin{proof}
Empty color classes or a color pair with no edge give an immediate no-instance. Otherwise number vertices by distinct codes in $\{1,\ldots,n\}$. Create one group for each color and one for each unordered color pair. Use one digit for every group and two incidence digits for every color pair, in base $H>(2g+1)(n+1)$. The target has digit one at every group position and digit $n$ at every incidence position.

A vertex of color $i$ supplies digit one at its color-group position and its code at every incidence position involving that color. An edge joining codes $u$ and $v$ supplies digit one at its pair-group position and digits $n-u$ and $n-v$ at the two corresponding incidence positions. All unspecified digits are zero. Each encoded integer is positive because of its group digit. A selection of at most one integer per group has no carries: its digit sums are at most $gn<H$. Equality at the group digits forces every group to be represented. Equality at each incidence digit forces the chosen edge endpoint to equal the chosen vertex. Thus the selection is precisely a clique and its edges. Conversely such a clique supplies the required sum. There are $O(r^2)$ digits, each using $O(\log(rn))$ bits, and at most one encoded integer per vertex or edge. The construction is a parameterized reduction. Multicolored clique is $\mathrm{W}[1]$-hard: from ordinary $r$-clique make $r$ color copies of every vertex and connect distinct copies exactly when their original vertices are adjacent; clique hardness follows by complementing the independent-set completeness result of \citet{DowneyFellows1995b}.
\end{proof}

\begin{proof}[Proof of Theorem~\ref{thm:wone59}]
Apply Lemma~\ref{lem:groupsum59}. A target outside $[0,\sum_i\max A_i]$ is an immediate no-instance and can be mapped to the fixed feasible renewal instance with one zero command, cap and promise zero, and threshold one. Otherwise set
\[
 M=\max_i\max A_i,\qquad T=2(g+1)M,\qquad K=(2g+1)T,
 \qquad \ell_i=\frac{i}{g+1},\quad i=1,\ldots,g.
\]
There are $2g+1$ equally likely histories: one with cap and ceiling zero, and, for each group $i$, two with caps and ceilings $\ell_i$ and $\ell_i+\max A_i/T$. The catalog contains zero, every baseline $\ell_i$, and an optional command $\ell_i+w/T$ for each distinct $w\in A_i$. All points lie in $[0,1]$. Options within group $i$ lie strictly between its baseline and the next group's baseline; the last option is strictly below one. Equal integers within a group can be merged without changing the sum problem. Charge zero for zero and the baselines, and charge $w/(2K)$ for the option representing $w$. Set
\begin{equation}\label{eq:budgetreduction59}
 m=2g+1,\quad D_0=\frac{2}{2g+1}\sum_i\ell_i,\quad
 B=D_0+\frac{W}{K},\quad \zeta=D_0+\frac{W}{2K}.
\end{equation}
The $2g+1$ caps are distinct, so $q_b=m$. The promise is feasible because
$\bar b=D_0+\sum_i\max A_i/K\geq B$ and zero is in the catalog.

Any feasible book must contain zero, since the zero-cap history has positive probability. Write $S(c)$ for the sum of the integers represented by its installed optional commands. With linear reward and free service, its exact gross value is
\[
 \min\{B,D(c)\},\qquad D(c)=\sum_j\pi_j\max(c\cap[0,b_j]).
\]
If $B\leq D(c)$, lotteries between zero and each last eligible command attain the required aggregate terminal mean. If $B>D(c)$, use those last commands surely and fill the remaining promise with deterministic service; the available service capacity is $\bar b-D(c)$. These constructions obey the original realized ceilings, not only the expected caps.

For an upper bound only, augment the book by all baselines, ignoring the augmented count. Its terminal capacity is
\[
 D(c)\leq D_0+\frac{1}{K}\sum_i\max\bigl(\{w:\ell_i+w/T\in c\}\cup\{0\}\bigr)
 \leq D_0+\frac{S(c)}{K}.
\]
The actual book, not the augmented comparison, pays charge $S(c)/(2K)$. Hence
\begin{equation}\label{eq:keyparam59}
 V_c(B)\leq \min\{B,D_0+S(c)/K\}-S(c)/(2K)
       =\zeta-\frac{|S(c)-W|}{2K}.
\end{equation}
Suppose its value reaches $\zeta$. Then $S(c)=W$ and both capacity inequalities must be equalities. All option weights are positive, so the second equality allows at most one optional command per group. The first equality forces every baseline into the actual book: if $\ell_i$ is absent, the dedicated history with cap $\ell_i$ has a strictly smaller last command, and adding $\ell_i$ strictly increases its positive-weight contribution to $D(c)$. Zero and the $g$ baselines therefore consume $g+1$ slots. Lemma~\ref{lem:groupsum59} now forces one option from each group. These $2g+1$ installed commands are within the original budget and encode a clique.

Conversely, choose the $g$ options representing a clique together with zero and all baselines. This book uses exactly $m$ commands, has $D(c)=D_0+W/K=B$, and has net value $\zeta$. This proves equivalence. The rational construction has polynomial bit length, and $m+q_b=4g+2=O(r^2)$, proving the parameterized hardness statement.
\end{proof}

This theorem does not assert $\mathrm{W}[1]$ membership, a sharp exponent lower bound, strong NP-hardness, or hardness under uniform positive charges. It establishes a parameterized barrier for the original restricted model, rather than for a larger abstract path problem. Its fine rational scales also explain why it is consistent with the additive algorithm and the pseudo-polynomial lattice algorithm. Exact enumeration is XP, whereas an exponent independent of both menu size and cap count would collapse the stated parameterized classes. The full earlier SUBSET SUM construction and the tight cap-count proof are retained in the electronic companion.
\section{A Tariff-Sensitive Tractability Frontier}\label{sec:tariff60}
The hardness result uses command-specific opening charges. This section isolates a positive counterpart without bounding the command allowance, the number of distinct caps, or the numerical denominators. Throughout the section, the common reward is $f(x)=rx$ with $r>0$ and service is free, $h_j=0$. Expected caps, realization ceilings, probabilities and catalog spacing remain arbitrary subject to Section~\ref{sec:model52}.

\begin{lemma}[Terminal capacity of a book]\label{lem:capacity60}
For a feasible book $c=(u_1<\cdots<u_s)$, put
\[
 D(c)=\sum_j\pi_j\min\{b_j,d_j(c)\},\qquad
 d_j(c)=\max\{u\in c:u\leq\tau_j\}.
\]
Its exact gross value at the accepted promise $B$ is $r\min\{B,D(c)\}$. Moreover,
\begin{equation}\label{eq:gain60}
 D(c)=u_1+\sum_{i=1}^{s-1}G_{u_i u_{i+1}},\qquad
 G_{uv}=\sum_{j:\tau_j\geq v}\pi_j[\min\{b_j,v\}-u]_+.
\end{equation}
\end{lemma}
\begin{proof}
The aggregate terminal mean cannot exceed either $B$ or $D(c)$. If $B\leq D(c)$, distribute $B-u_1$ within the individual intervals $[u_1,\min\{b_j,d_j(c)\}]$, using their probability weights. Adjacent eligible commands implement each resulting terminal mean with zero service. If $B>D(c)$, first attain every terminal capacity. Histories with $b_j>d_j(c)$ then have a deterministic terminal command $d_j(c)$ and residual service capacity $b_j-d_j(c)$. These capacities total $\bar b-D(c)\geq B-D(c)$. Filling them meets the promise exactly and respects the realization ceiling because the final deterministic obligation is at most $b_j\leq\tau_j$. This attains the upper bound. For the identity, fix a history and telescope its eligible prefix after truncation at $b_j$. Every increment is precisely the summand in \eqref{eq:gain60}; the initial contribution is $u_1$ since $u_1\leq\min_j b_j$. Sum with weights $\pi_j$.
\end{proof}

Choose a standard opening fee $\rho_0\geq0$ and let $E=\{i:\rho_i\ne\rho_0\}$ and $d=|E|$. This parameter counts exceptional \emph{commands}, not distinct fee values. A most frequent catalog fee minimizes $d$ and is found without optimizing any allocation. Equal-frequency modes are resolved by the smallest rational fee; any mode gives the same minimum exception count. For $J\subseteq E$, require exactly those exceptional commands in $J$ and forbid $E\setminus J$. Denote the allowed indices by $A_J$. An edge $i\to j$ with $i<j$ is admissible when $i,j\in A_J$ and no required index lies strictly between them. A starting index must precede or equal the first required index; a final index must follow or equal the last required index. These conventions are vacuous when $J$ is empty.

\begin{theorem}[Exact allocation with few tariff exceptions]\label{thm:tariff60}
Under the assumptions of this section, exact optimization and original-policy recovery require
\begin{equation}\label{eq:tariffcost60}
 O(kN^2+2^d mN^2)
\end{equation}
rational arithmetic operations, after replacing $m$ by $\min\{m,N\}$. The algorithm has bit complexity $2^d\operatorname{poly}(L)$ for binary input length $L$. It uses no discretization of the promise or capacities. In particular, a uniform nonnegative opening fee gives a polynomial-time algorithm with a strongly polynomial arithmetic count, for unrestricted $m$ and $q_b$.
\end{theorem}
\begin{proof}
Precompute all gains in \eqref{eq:gain60}. For each $J$, let $T^J_s(j)$ be the largest capacity of an admissible $s$-command prefix ending at $a_j$. Use $-\infty$ for absent states and initialize
\[
 T^J_1(j)=a_j
 \quad\text{if }j\in A_J,\quad a_j\leq\min\{B,\min_i b_i\},
 \quad j\leq\min J,
\]
where the last condition is omitted for $J=\varnothing$. The recurrence is
\begin{equation}\label{eq:tariffdp60}
 T^J_s(j)=\max_{i<j:\,i\to j\text{ admissible}}
       \{T^J_{s-1}(i)+G_{a_i a_j}\},\qquad 2\leq s\leq m.
\end{equation}
Accept only final indices $j\geq\max J$, with the condition again vacuous for an empty set. A valid start, edges that cannot skip a required index, and a valid end together force every member of $J$ to appear. Conversely, every book with exceptional set $J$ follows exactly such a path. For fixed $J$ and $s$, every accepted book pays the same charge
\[
 K^J_s=(s-|J|)\rho_0+\sum_{i\in J}\rho_i.
\]
By Lemma~\ref{lem:capacity60}, its value is nondecreasing in capacity. Keeping only the maximum capacity at each state therefore loses no optimal book. Maximize $r\min\{B,T^J_s(j)\}-K^J_s$ over accepted states and all subsets $J$. Predecessors recover an attaining book, and the lemma recovers a feasible original allocation. States with $s<|J|$ cannot reach a valid end and contribute nothing.

Computing the gains costs $O(kN^2)$. Prefix counts of required indices test edge admissibility in constant time, so each subset costs $O(mN^2)$. Capacity entries are sums of at most $N$ gains; equivalently, every attained entry is a sum of $k$ weighted truncated commands. A common product of the input denominators bounds their encoding length polynomially in $L$. Exact comparisons, the value calculation and the interval-filling recovery consequently have polynomial bit cost per operation. Neither a common denominator's numerical value nor a grid size enters \eqref{eq:tariffcost60}. This proves fixed-parameter tractability in $d$ and the uniform-fee conclusion.
\end{proof}

The distinction from Theorem~\ref{thm:wone59} is substantive. Its linear/free-service hard instances remain within this section's physical primitives but have command-specific fees that need not have few exceptions. Thus a cap-count and command-count lower bound coexists with an exception-count upper bound. Enumerating exceptions costs an exponential function of $d$ multiplying a polynomial whose exponent is independent of $d$; enumerating all books costs a parameter-dependent exponent in $N$. Uniform fees with nonlinear rewards or positive service costs are outside Theorem~\ref{thm:tariff60}, not instances silently solved by it.

\begin{corollary}[The uniform-fee decision frontier]\label{cor:fee60}
Fix the physical input and $B$. Let $D_s$ be the largest terminal capacity among feasible books with exactly $s$ commands, omitting infeasible sizes. Under a uniform fee $\rho$, the value is
\[
 V(\rho)=\max_{1\leq s\leq m}\{r\min(B,D_s)-s\rho\}.
\]
It is convex, nonincreasing and piecewise affine in $\rho$. The smallest optimal book size is nonincreasing in $\rho$, and every switching fee is an intersection of two displayed affine functions. The capacity tables compute the full fee frontier without rerunning the resource allocation problem.
\end{corollary}
\begin{proof}
The formula follows from the theorem. A maximum of decreasing affine functions is convex and nonincreasing. If $s_1,s_2$ are optimal at $\rho_1<\rho_2$, their two optimality inequalities imply $(s_1-s_2)(\rho_2-\rho_1)\geq0$. Pairwise intersections give all possible switches, with dominated intersections discarded. No diminishing-returns property in the command budget is assumed.
\end{proof}
\section{Near-Standard Tariffs and Implementable Decision Frontiers}\label{sec:robust63}
An exact standard tariff is an institutional restriction, not a rounding convention. A common installation charge may represent a standardized configuration procedure, while a few commands require separate certification or integration. The count of exceptional commands is observable before allocation. Small deviations on many commands, however, can make that exact count large. This section shows how to approximate such fees while retaining a certificate for the \emph{original} tariff and the original physical policy. The physical assumptions of Theorem~\ref{thm:tariff60} are needed for its tractable inner optimization; the perturbation bound itself does not require linear rewards or free service.

\subsection{A fee-perturbation certificate}
Fix the physical primitives and the feasible book family $\mathcal C_B$. Write $W_c(B)$ for the gross value of a book, and $V_\rho=\max_{c\in\mathcal C_B}\{W_c(B)-\sum_{i\in c}\rho_i\}$. Replace the fee vector $\rho$ by a nonnegative vector $\widetilde\rho$, without changing any other input. Put $\delta_i=\rho_i-\widetilde\rho_i$, and, with $\bar m=\min\{m,N\}$, define
\[
 E_+=\sum_{\ell=1}^{\bar m}(\delta_+)_{[\ell]},\qquad
 E_-=\sum_{\ell=1}^{\bar m}((-\delta)_+)_{[\ell]},
\]
where descending order statistics are padded with zeros. These quantities depend only on the two tariffs and the command allowance.
\begin{theorem}[Certified transfer to the original tariff]\label{thm:robust63}
Suppose $\widetilde c$ and an original-feasible policy attain $V_{\widetilde\rho}$. Evaluating that same policy under $\rho$ gives the certified interval
\begin{equation}\label{eq:robust63}
 L=V_{\widetilde\rho}-\sum_{i\in\widetilde c}\delta_i
 \ \leq\ V_\rho\ \leq\ U=V_{\widetilde\rho}+E_-.
\end{equation}
Consequently,
\[
 0\leq U-L\leq E_++E_-\leq\|\rho-\widetilde\rho\|_1
 \leq N\|\rho-\widetilde\rho\|_\infty,
 \qquad E_++E_-\leq2\bar m\|\rho-\widetilde\rho\|_\infty.
\]
The policy preserves the book allowance, every history's constraints, and the accepted aggregate promise exactly. If the surrogate optimum exceeds every other book's surrogate value by more than $E_++E_-$, its book remains uniquely optimal under $\rho$.
\end{theorem}
\begin{proof}
For every admissible book $c$, at most $\bar m$ fee differences are summed, so
$-E_-\leq\sum_{i\in c}\delta_i\leq E_+$. Its original net value equals its surrogate net value minus that sum. Maximizing gives the upper bound in \eqref{eq:robust63}; retaining the attaining policy gives the lower bound. Subtracting the bounds proves the regret statement. Positive and negative differences occupy disjoint coordinates, which proves the $\ell_1$ inequality; each one-sided sum has at most $\bar m$ terms, proving the second norm bound. Feasibility involves no opening fee, so the same allocation remains feasible. Finally, perturbation can decrease the difference between the chosen book and any competitor by at most $E_++E_-$. A strictly larger surrogate separation therefore preserves its strict ordering.
\end{proof}

Thus a uniform approximation within $\eta$ yields an additive guarantee $2\bar m\eta$, and one-sided fee changes yield $\bar m\eta$. More informative book-specific bounds follow directly from \eqref{eq:robust63}. A numerically small fee dispersion does not make arbitrary fees \emph{exactly} fixed-parameter tractable. It creates a certified approximation neighborhood of the exact theorem, with an explicit price for simplifying the tariff.

\subsection{Selecting a sparse-exception surrogate}
For an integer $0\leq d<N$, consider surrogates that leave at most $d$ original fees unchanged as exceptions and replace all remaining fees by one nonnegative standard. Let $n_0=N-d$, and sort the original fees as $r_1\leq\cdots\leq r_N$, breaking equal-fee ties by catalog index. For each consecutive window of $n_0$ observations, assign its lower median as the standard and leave all observations outside the window unchanged.
\begin{proposition}[Minimum absolute tariff distortion]\label{prop:projection63}
A window of minimum absolute deviation produces a surrogate minimizing $\|\rho-\widetilde\rho\|_1$ among all surrogates with at most $d$ unchanged exceptions. Its actual number of fees different from the chosen standard is at most $d$. The minimum distortion is nonincreasing in $d$. Hence, for any tolerance $\varepsilon>0$, binary search over $d\in\{0,\ldots,N-1\}$ finds the smallest exception allowance in this surrogate class with distortion at most $\varepsilon$.
\end{proposition}
\begin{proof}
For a fixed standard, an optimal choice of replaced entries consists of $n_0$ entries closest to that standard. In a sorted list there is such a choice forming a consecutive window: if an omitted entry lies between two chosen entries, its distance does not exceed the larger endpoint distance, so exchanging that endpoint cannot increase the objective. Repeating removes all holes, with equal distances resolved consistently. For a fixed window, a median minimizes the sum of absolute deviations; its lower median is nonnegative. Minimizing over all windows therefore attains the global minimum. Entries outside the window are the only possible exceptions, and an outside fee equal to the median further reduces their actual count. Enlarging the exception allowance retains every previously feasible surrogate, proving monotonicity. At $d=N-1$ the distortion is zero, so the stated binary search terminates.
\end{proof}

Sorted prefix sums evaluate all window deviations in linear time after sorting. The retained implementation re-sorts at each binary-search step and therefore uses $O(N\log^2 N)$ comparison/arithmetic preprocessing; sorting once reduces this to $O(N\log N)$. We report the implemented, rather than the improved hypothetical, preprocessing bound. Applying Theorem~\ref{thm:tariff60} to the selected surrogate takes $2^d\operatorname{poly}(L)$ bit operations and returns the original-policy interval \eqref{eq:robust63}. The selected $d$ minimizes the stated $\ell_1$ distortion criterion, not the realized regret or wall-clock cost. The two one-sided order-statistic budgets can certify a smaller error than the projection objective. A declared computational guard may still stop a large-$d$ request.

\subsection{Exact fees, installed commands, and policy changes}\label{sec:frontier63}
Consider four equally likely histories with caps and ceilings $(0,1/3,2/3,1)$, catalog $(0,1/3,2/3,1)$, $m=4$, $B=2/5$, reward $f(x)=x$, and free service. At a uniform opening fee $\rho$, the relevant capacity envelopes give
\[
 V(\rho)=\max\{2/5-3\rho,\;1/3-2\rho,\;-\rho\}.
\]
The four-command line is dominated under the smallest-size tie convention. For $0\leq\rho<1/15$, an optimal representative is $(0,1/3,2/3)$; for $1/15\leq\rho<1/3$, it is $(0,2/3)$; and for $\rho\geq1/3$, it is $(0)$. At a switching fee the smaller book is selected. The first book has terminal capacity $5/12$ and delivers the entire promise through terminal lotteries. The second has capacity $1/3$ and supplies aggregate preliminary service $1/15$. The singleton supplies all $2/5$ through preliminary service. Thus the frontier specifies both the installed commands and the allocation mechanism, not only the envelope value. These are deterministic representatives; uniqueness is not asserted at ties.

Now take the nonuniform fees $(67/2000,33/1000,203/6000,199/6000)$ and tolerance $1/1000$. The projection selects one exceptional command and returns the original book $(0,1/3,1)$. Its original net value is $901/3000$, and the independently checked upper bound is $451/1500$, giving gap $1/3000$. Here $E_+=1/3000$ and $E_-=1/6000$. The change in book composition is not described by a monotonic nesting claim: only the smallest optimal \emph{size} is monotone along a uniform-fee frontier. The code/data package also supplies complete rational frontiers and reconstructed policies for catalogs of 17, 65, and 129 commands.
\section{Merged-Origin Resource-Deficit Paths}\label{sec:deficit56}
The capacity potential does more than guarantee repair after rounding. It identifies a resource coordinate in which every feasible anchor has the same initial state and every feasible book has the same terminal state. This removes the extra anchor loop in the resource algorithm. The construction applies to the heterogeneous kernels of Theorem~\ref{thm:packing54}.

\subsection{An exact source-independent coordinate}
Write
\[
 p(u)=\sum_j\pi_j\min\{b_j,u\},\qquad p(\dagger)=\bar b,
 \qquad C_{uv}=T_{uv}+Q_{uv}=p(v)-p(u).
\]
For the terminal edge, $C_{u\dagger}=\bar b-p(u)$. Every feasible anchor $a\leq\min\{B,\min_jb_j\}$ satisfies $p(a)=a$. If edge $e$ is assigned resource $x_e$, define its unused capacity $w_e=C_e-x_e$ and its deficit payoff $\Psi_e(w)=\Phi_e(C_e-w)$. Opening charges are kept separate from $\Phi_e$ throughout. Let $R=\bar b-B$.

\begin{theorem}[Merged-origin identity]\label{thm:deficit56}
Under the common ordered catalog, prefix eligibility, and increasing concave history-specific rewards of Theorem~\ref{thm:packing54}, the original value equals
\begin{equation}\label{eq:deficit56}
 \max_{\substack{c:\,1\leq |c|\leq m\\a=\min c\leq\min(B,\min_j b_j)}}
 \left\{\sum_j\pi_j f_j(a)-\sum_{z\in c}\rho(z)
 +\max_{\substack{0\leq w_e\leq C_e\\\sum_{e\in P(c)}w_e=R}}
           \sum_{e\in P(c)}\Psi_e(w_e)\right\}.
\end{equation}
All anchors start at cumulative deficit zero and end at cumulative deficit $R$. Every feasible deficit vector on a chosen book yields a policy at the original promise, within the original command budget, with at least its displayed value. No anchor identifier is needed in a Bellman state once the current command and number of selected commands are known.
\end{theorem}
\begin{proof}
For a fixed book with anchor $a$, capacity conservation gives $\sum_e C_e=\bar b-a$. Hence
\[
 \sum_e x_e=B-a\quad\Longleftrightarrow\quad
 \sum_e(C_e-x_e)=\bar b-B=R.
\]
The transformation is a bijection between the two resource boxes intersected with their equalities. Apply the exact component-packing identity to obtain \eqref{eq:deficit56} and its recovery claim. Along a prefix ending at $u$, the cumulative deficit is
\[
 d=p(u)-a-\sum_{e\text{ in the prefix}}x_e.
\]
At the anchor it is zero; traversing $u\to v$ adds $C_{uv}-x_{uv}$. Continuation payoffs and transitions depend only on $u$, the remaining command allowance, and $d$, not on the earlier anchor. A Bellman merge retains the best accumulated value among such prefixes. The recovered full path retains its own anchor through predecessor pointers. Partial paths need not themselves encode a canonical history policy: the full-path packing theorem supplies implementation after the equality is restored.
\end{proof}

The same argument gives an abstract reuse principle. On an acyclic graph, suppose every resource capacity has the form $C_{uv}=p(v)-p(u)\geq0$, all paths share a sink, and a source $a$ requires resource $B-p(a)$ with $p(a)\leq B$. Complementing edge resources makes the total unused capacity $p(\dagger)-B$ independent of the source. Any additive path objective and edge-count bound then admit the same source merge. Concavity is needed for the approximation and fast convolution below, not for this algebraic identity. The renewal-specific content is the exact derivation of those capacities and payoffs from individual expected and realized constraints; an arbitrary constrained path need not have this potential structure.

\subsection{Centered approximation and exact repair}
For rational quadratic primitives put
\[
 r_{\max}=\max_jr_j,\qquad G=\max_j\gamma_j,
 \qquad \lambda_0=(r_{\max}-G)/2,\qquad K=(r_{\max}+G)/2>0.
\]
Every supergradient of $\Phi_e$ lies between $-G$ and $r_{\max}$ on its interval. Consequently
$g_e(w)=\Psi_e(w)+\lambda_0w$ is concave and $K$-Lipschitz. Adding this common linear term changes every exactly feasible path value by the same constant $\lambda_0R$; it does not change the optimizer.

For $\varepsilon>0$ and $R>0$, set
\begin{equation}\label{eq:mesh56}
 \eta=\frac{\varepsilon}{2Km},\qquad
 Q=\left\lfloor\frac{R}{\eta}\right\rfloor+1,\qquad
 \mathcal I=\{s\in\{0,\ldots,Q-1\}:R-m\eta\leq s\eta\leq R\}.
\end{equation}
Let $D_{\ell,v}(s)$ denote the best accumulated centered value at command $a_v$, with exactly $\ell$ selected commands and cumulative deficit $s\eta$. Unreachable entries are $-\infty$. All feasible anchors are initialized simultaneously:
\begin{align}
 D_{1,i}(0)&=\sum_j\pi_jf_j(a_i)-\rho_i,
       &&a_i\leq\min\{B,\min_jb_j\},\label{eq:init56}\\
 D_{\ell,v}(s)&=\max_{\substack{u<v,\ 0\leq t\leq s\\t\eta\leq C_{uv}}}
 \{D_{\ell-1,u}(s-t)+g_{uv}(t\eta)-\rho_v\},\label{eq:DP56}\\
 S_\eta&=\max_{\substack{1\leq\ell\leq m,\ u,\ s\in\mathcal I\\0\leq t\leq s,\ t\eta\leq C_{u\dagger}}}
 \{D_{\ell,u}(s-t)+g_{u\dagger}(t\eta)\}.\label{eq:terminal56}
\end{align}
The tail does not install a new command. A book with $s$ selected commands has $s$ resource edges including that tail.

\begin{theorem}[Centered deficit approximation]\label{thm:approx56}
Equations~\eqref{eq:init56}--\eqref{eq:terminal56} and exact fixed-book recovery produce an original-instance interval
\begin{equation}\label{eq:interval56}
 \underline V\leq V^*(B)\leq S_\eta-\lambda_0R+Km\eta,
 \qquad S_\eta-\lambda_0R+Km\eta-\underline V\leq\varepsilon.
\end{equation}
The guarantee allows heterogeneous rewards, arbitrary full-catalog eligibility classes, arbitrary nonnegative charges, and the original budget. For $R=0$ the same recurrence uses only deficit zero and is exact, including with nonzero service costs.
\end{theorem}
\begin{proof}
Round each deficit on an optimal continuous path down to a multiple of $\eta$. Its cumulative rounded deficit is no greater than $R$ and differs from $R$ by less than $m\eta$. Every prefix stays in the grid. Lipschitz continuity loses at most $Km\eta$ from the centered optimal value $V^*(B)+\lambda_0R$. Thus $S_\eta\geq V^*(B)+\lambda_0R-Km\eta$, proving the upper bound.

Conversely, recover a maximizing grid path. Increase its deficits, without exceeding their individual capacities, until their sum is $R$. This is always possible because the path's total capacity is $\bar b-a\geq\bar b-B=R$. The total increase is at most $m\eta$. The centered payoff decreases by at most $Km\eta$; after subtracting $\lambda_0R$, the uncentered path is worth at least $S_\eta-\lambda_0R-Km\eta$. Component packing gives an original feasible policy with no smaller value. Exact fixed-book allocation may improve it further, so the resulting feasible lower value satisfies the same bound. This proves \eqref{eq:interval56}. When $R=0$, nonnegative deficits must all be zero, so there is no rounding, repair, or approximation error.
\end{proof}

\subsection{Convolution with holes and bit complexity}
The merge can create unreachable gaps in a Bellman row. Concavity of the row itself is neither assumed nor generally true.

\begin{lemma}[Concave convolution with arbitrary finite support]\label{lem:holes56}
Let $A(z)\in\mathbb R\cup\{-\infty\}$ be arbitrary on $0,\ldots,Q-1$, and let $g(t)$ be a finite concave sequence on $0,\ldots,L$. The smallest maximizing input index in
\[
 M(s)=\max_{\max(0,s-L)\leq z\leq s}\{A(z)+g(s-z)\}
\]
is nondecreasing across rows having at least one finite feasible column. All values and predecessors can be computed in $O(Q\log(Q+1))$ comparisons and additions, including arbitrary holes in the support of $A$.
\end{lemma}
\begin{proof}
Delete columns with $A(z)=-\infty$ and rows whose feasible interval contains no remaining column. Concavity gives the anti-Monge inequality
\[
 [A(z)+g(s-z)]+[A(z')+g(s'-z')]
 \geq[A(z')+g(s-z')]+[A(z)+g(s'-z)]
\]
whenever $s<s'$, $z<z'$, and the four entries are feasible. If two maximizing indices crossed in the wrong order, their crossed rectangle would be feasible because $0\leq s-z\leq L$ is a staircase band. The inequality and the smallest-index tie rule would contradict that crossing. The maximizing indices are therefore monotone. Sorted finite columns permit feasible-row detection by interval searches. Divide-and-conquer row maximization scans a total $O(Q)$ candidate positions at each of $O(\log(Q+1))$ levels, with shared endpoint candidates included. Empty rows are discarded before recursion; they are not allowed to impose an artificial predecessor bound on another row.
\end{proof}

\begin{proposition}[Arithmetic, storage, and verification]\label{prop:cost56}
For rational quadratic inputs the centered deficit algorithm has arithmetic bound
\begin{equation}\label{eq:cost56}
 O\bigl(kN^2Q\log(k+1)+mN^2Q\log(Q+1)\bigr).
\end{equation}
Its cached implementation stores $O((N^2+mN)Q+kN^2)$ rational entries and has bit complexity polynomial in binary input length and numerical $Q$. The extra factor $N$ for separately enumerated anchors in \eqref{eq:fallback54} is absent. An independent checker can reconstruct all kernels and verify Bellman upper-potential inequalities in
$O(k^2N^2Q+mN^2Q^2)$ rational operations, without using the optimizer's monotone search.
\end{proposition}
\begin{proof}
There are $O(N^2)$ ordinary and terminal edges. Sorted weighted terminal segments and sorted service saturation thresholds give piecewise-linear terminal profiles and piecewise-quadratic service profiles. Their endpoints and cumulative values are rational. Direct evaluation by capped allocation already gives the first term in \eqref{eq:cost56}; caching endpoints and binary searching them can reduce its constant and history dependence. For each selected count and edge, Lemma~\ref{lem:holes56} gives the second term. The predecessor table has $mNQ$ entries. Caching every edge grid adds $N^2Q$ entries; retaining all component descriptions adds at most $kN^2$. Thus the storage claim includes the edge cache rather than counting Bellman states alone.

A kernel value involves finite sums and products of input rationals, a grid rational, and at most one capped-service water-level division. Sums of at most $k$ rational reciprocals have encoding length polynomial in the total input encoding. A path combines at most $m$ such values. Therefore all represented integers have bit length polynomial in the input length and $\log(Q+1)$; multiplying arithmetic counts by polynomial integer-arithmetic costs gives the stated bit bound. A common denominator may have exponentially large numerical value, which is not a polynomial-time claim in its binary encoding.

For independent verification, supplied Bellman values must dominate every feasible anchor and every predecessor-plus-edge transition. No transition into a claimed unreachable entry may be omitted. Every accepted terminal combination must be covered by the reported terminal bound. These inequalities inductively bound every grid path; an attaining path certifies the grid optimum used in the repair calculation. A separately checked original lottery supplies the lower value. Direct service-threshold testing costs at most quadratic time in histories per kernel value; explicitly examining every predecessor deficit gives the stated conservative checker bound. The checker uses neither the forward optimizer nor its divide-and-conquer routine.
\end{proof}

With $\delta=\varepsilon/K$ and $R\leq1$, substitution gives a budget-quadratic, catalog-quadratic term of order $m^2N^2\delta^{-1}\log(m/\delta+2)$ rather than its catalog-cubic predecessor. The guarantee is additive in signed net value. It is not polynomial in $\log(1/\varepsilon)$, and a successful optimization need not imply a cheap independent verification; both costs are measured separately.
\section{When a Price Bound Is Already an Exact Decision}\label{sec:pricebridge60}
A useful algorithmic distinction is not simply between an integer and a continuous method, but between a price bound that supports one implementable book and a bound that mixes books. Fix the original feasible book family $\mathcal C_B$. For $c\in\mathcal C_B$, let $F_c(t)$ denote its exact net value at aggregate target $t\in[c_1,\bar b]$. For a price $\lambda$, define
\[
 H_c(\lambda)=\max_t\{F_c(t)-\lambda t\},\qquad
 U(\lambda)=\lambda B+\max_{c\in\mathcal C_B}H_c(\lambda).
\]
A book is price-active if it attains the second maximum. The maximizing target set $I_c(\lambda)=\argmax_t\{F_c(t)-\lambda t\}$ is a closed interval by concavity. Its endpoints can be recovered by summing the smallest and largest individual maximizing targets. The exact price-path oracle in the companion computes $U(\lambda)$ without enumerating all books.

\begin{theorem}[Same-book support and a distance bound]\label{thm:pricebridge60}
For any price $\lambda$, $U(\lambda)=V^*(B)$ if and only if a price-active book $c$ has $B\in I_c(\lambda)$. For any price-active $c$, if its fixed-book value is $L$-Lipschitz, then
\begin{equation}\label{eq:pricedistance60}
 0\leq U(\lambda)-F_c(B)
 \leq (L+|\lambda|)\operatorname{dist}(B,I_c(\lambda)).
\end{equation}
For the quadratic primitives, the common bound in \eqref{eq:L52} is sufficient. With free service, common linear reward and zero opening fees, a single exact price certificate always exists: use $\lambda=r$ when $B\leq\max_c D(c)$, and $\lambda=0$ otherwise.
\end{theorem}
\begin{proof}
If $B\in I_c(\lambda)$ for an active book, then $U(\lambda)=F_c(B)\leq V^*(B)$, and weak duality gives equality. Conversely, take an optimal original book. Equality of the global bound with its value forces equality both in its individual price bound and in the maximization over books; it is active and its maximizing target interval contains $B$. For \eqref{eq:pricedistance60}, project $B$ onto $I_c(\lambda)$ at $T$. Then
$U(\lambda)-F_c(B)=F_c(T)-F_c(B)+\lambda(B-T)$, which gives the stated bound. A fixed-book allocation can be moved to a higher or lower aggregate target by filling or emptying residual individual target intervals. The total weighted change is exactly the aggregate change, so individual $L$-Lipschitz responses imply the same bound for $F_c$.

In the zero-fee linear case, $F_c(t)=r\min\{t,D(c)\}$ by Lemma~\ref{lem:capacity60}. At price $r$, every feasible book is active and its maximizing interval is $[c_1,D(c)]$. Choose a capacity-maximizing book when $B$ is below that capacity. At price zero, capacity-maximizing books are active on $[D(c),\bar b]$. This covers the remaining case.
\end{proof}

The interval must belong to \emph{one} active book. Merely placing $B$ between targets attained by different active books certifies a relaxation, not an original policy. Even a uniform positive fee can create this distinction. Take two equally likely histories with $b=\tau=(0,1)$, catalog $(0,1)$, $m=2$, common reward $f(x)=x$, free service, and fee $1/4$ for each installed command. At $B=1/8$, the one-command book has net value $-1/4$, whereas the two-command book has value $-3/8$. At $\lambda=1/2$, their price intercepts coincide at $-1/4$, producing the upper bound $-3/16$ and gap $1/16$. The tariff algorithm nonetheless solves this instance exactly. Thus its positive result is not merely a restatement of zero Lagrangian gap.

These certificates suggest observable diagnostics, not a universally validated routing rule. The implemented hybrid and its unfavorable as well as favorable outcomes are evaluated in Section~\ref{sec:study59}. The exact tariff theorem, a same-book support witness, and a resource approximation guarantee have different scopes and different checking costs.
\section{Computational Evidence and the Cost of Certification}\label{sec:study59}
The experiments address three separate questions: how an exact exception-count algorithm scales; whether an inexpensive price bound already supports an implementable book; and how much end-to-end certification costs beyond optimization. All inputs are constructed and disclosed. The study does not establish adoption, calibration to a provider, or a universally superior routing rule. Tables are collected after the references; complete input, source, execution, and proof records accompany the paper.

\subsection{Preserved evidence and an expanded frozen design}
The earlier study comprises 90 distinct model specifications, 270 budget-tagged cases, and 1,260 method requests. In its three-second standard-tariff panel, the tariff, price, deficit, and enumeration methods attain an independently checked rational target in 30, 29, 12, and 19 of 30 requests, respectively. SCIP attains its numerical target in 30 of 30. The corresponding exceptional-tariff counts are the same. In the curved panel, the counts are 28 for price, 10 for deficit, and 18 for enumeration; the tariff method is mathematically inapplicable, and SCIP reports 27 numerical successes. These results, the preceding resource-grid experiments, and all unfavorable outcomes are preserved unchanged. A favorable guarantee is not interpreted as evidence that the earlier deficit implementation was fastest.

The expanded design contains 44 distinct physical/tariff specifications, 88 budget-tagged cases, and 405 declared requests. Its controlled scaling panel crosses $(N,k,m)=(17,12,8),(65,32,16),(129,48,32)$ with $d\in\{0,1,2,4,6,8,10,12\}$ and one- and six-second total allowances. Each cell runs tariff, price, enumeration, SCIP, and a specified hybrid. Eight further tariff requests use 30 seconds on the 17-command cases. Three 13-exception requests test the engineering guard. A separate phase-split panel varies the optimization allocation between 35\% and 80\% of a three-second allowance. Five numerical-encoding cases use probability constructions based on 16, 64, 256, 1,024, and 4,096 bits, for tariff and price under 15 seconds. These are controlled cells, not random replicates.

A separately specified adversarial panel has twelve cases at $(N,k,m)=(19,10,6)$ and $(41,20,12)$. It varies exception density, a low accepted promise, many small fee deviations, positive service costs and curvature, eligibility geometry, and alternating fee levels. It does not reuse the earlier seeds or the hardness reduction. Its purpose is broader stress coverage, not a field-data claim. It adds the robust-tariff and deficit methods, giving 84 six-second requests. The full formulas and rational inputs are frozen before the expanded execution.

\subsection{The exception parameter in time, memory, and proof size}
Table~\ref{tab:scaling63} reports the controlled success counts with their full denominators. At six seconds the exact tariff method certifies all eight 17-command cases, five of eight 65-command cases, and three of eight 129-command cases. The largest certified exception counts at those sizes are 12, 6, and 2, respectively. Thus the parameterized result is reflected in an increasingly costly exception loop, not an assertion that every value up to the implementation guard is inexpensive. All eight long-allowance tariff requests certify the target.

Table~\ref{tab:d63} separates optimization, serialization, checking, resident memory, and uncompressed proof size for every exception count at the fixed 17-command scale. Increasing $d$ from zero to twelve changes the optimization time from 0.005 to 0.819 seconds and the independent check from 0.209 to 1.116 seconds; proof size increases from 4,143 to 3,964,865 bytes. At 65 commands and eight exceptions, an optimizer can finish while the parent deadline prevents completion of verification. That request remains a failure of the six-second end-to-end target. Proof files in the entire expanded panel reach 5,572,345 uncompressed bytes. Reported phase costs condition on the phase actually being reached, and optimizer and checker peak memory are distinct measurements.

The tariff algorithm certifies all five numerical-encoding cases, as does price, without a resource grid. These finite cases test exact rational arithmetic with large encodings; they do not remove bit-operation costs from the theorem. Enumeration attains none of the 48 controlled scaling targets within these allowances. The result illustrates the difference between the exception-count compression and book enumeration at these sizes; it is not an empirical proof of an asymptotic lower bound. The adversarial panel has rational successes of 4/12 for tariff, 7/12 for robust tariff, 12/12 for price, 8/12 for hybrid, and 0/12 for both enumeration and deficit. SCIP reports 12/12 numerical successes, with only original-feasible lower policies independently certified rationally.

\subsection{Same-book support: an independent exhaustive diagnostic}\label{sec:root63}
The expanded large-instance price logs contain 69 requests with an observed root probe, and none of the particular returned active books supports the promise at its recorded probe price. This does \emph{not} prove that the whole active-book family lacks support: ties and a nonminimizing probe price matter. We therefore add a separate, fully enumerated diagnostic rather than treat absence in a trace as a structural conclusion.

This diagnostic uses four equally likely caps $(0,1/3,2/3,1)$, seven catalog commands $i/6$, two ceiling patterns, allowances two and four, three promise fractions of the mean cap, and four fee patterns. The complete factorial design has 48 cases, frozen before its own execution. For every feasible book, an independent routine constructs the three vertices of its exact linear/free-service value function: its first command, terminal capacity, and mean cap. The upper concave hull of the union gives the exact minimum price bound at the promise. All active books and their maximizing intervals are then enumerated, independently of the price-path recurrence. Finally, the price implementation is run and its certificate is checked against the external input.

Across 1,176 enumerated feasible books, 20 cases have same-book support and 28 have a strictly positive minimum root gap; the largest gap is $1/20$. The equality characterization and the distance bound of Theorem~\ref{thm:pricebridge60} hold in every checked case. Table~\ref{tab:root63} contrasts observed branching and checking work by this independently established classification. All 48 final price certificates are exact. The comparison identifies a mechanism on this balanced small panel; it is not a causal estimate or a population frequency, and its local timings are not pooled with the earlier execution environment.

\subsection{A tested hybrid, and the limits of a routing claim}
The expanded hybrid spends at most one quarter of its optimization allowance on a one-node, three-price probe. If the certified gap is small enough it returns that proof. Otherwise, in the linear/free-service regime it tries the near-standard projection with an exception guard of four when the projected distortion meets the tolerance; failing that test, it restarts price search. The restart cost is charged rather than described as a free warm start. This rule was executed, not merely proposed.

Its 27/48 controlled scaling successes compare with 25/48 for tariff and 24/48 for price. At the largest size and six seconds it certifies three of eight cases while price certifies none; at one second it supplies no advantage at either larger size. In the adversarial panel price certifies all twelve cases and the hybrid eight. These outcomes do not validate a generally optimal method-selection hierarchy. They support reporting the observable inputs---physical applicability, exact or approximate exception count, and an available support witness---together with the full cost of checking the resulting decision.

The structural regression suite additionally checks 198 robust instances against 198 exhaustive fee projections and 6,623 book-error comparisons, rejects 108 corrupt certificates, and exercises exact exception counts through twelve and the guard at thirteen. The electronic companion explains the external-input serialization repair and the complete independent replay. Neither successful post-study replay nor a new packaging build retroactively converts an interrupted request into an on-time success.

\section{Conclusion}\label{sec:conclusion63}
The finite-catalog model admits an exact resource-path representation, but representation alone does not determine tractability. With common linear rewards and free service, arbitrary command-specific charges retain parameterized hardness in the command allowance and cap count, while a small number of exceptions to a standard charge permits an exact fixed-parameter algorithm. The uniform-fee frontier gives explicit installation decisions, and the perturbation certificate extends that algorithm to near-standard tariffs with a bound on the value lost by an original-feasible policy. The remaining conservation, approximation, heterogeneous-reward, prototype, and price results specify complementary regimes rather than a universal algorithm.

Computational certification has its own costs. Larger catalogs, larger exception counts, and independent checking can exhaust a deadline even when an optimizer produces a candidate. The expanded controlled and adversarial records, and the independent root-support diagnostic, make these distinctions observable. A practical implementation should preserve both the policy and its precise scope of guarantee, including the model input, fee approximation, resource equality, and time allowed for verification.

\paragraph{Code and data.} The accompanying archive contains the complete constructed inputs, unchanged historical execution sources and records, current independent checkers, reconstruction instructions, and all retained derivations. No proprietary operational data are used.
\clearpage\phantomsection\label{refs-start}
\begin{thebibliography}{99}\raggedright
\bibitem[Aggarwal et al.(1987)Aggarwal, Klawe, Moran, Shor, and Wilber]{Aggarwal1987}
Aggarwal A, Klawe MM, Moran S, Shor P, Wilber R (1987) Geometric applications of a matrix-searching algorithm. \emph{Algorithmica} 2:195--208. doi:10.1007/BF01840359.
\bibitem[Ahmadi et al.(2025)Ahmadi, Raith, and Jalili]{Ahmadi2025}
Ahmadi S, Raith A, Jalili M (2025) A fast and simple algorithm for the resource constrained shortest path problem. \emph{33rd Annual European Symposium on Algorithms}, LIPIcs 351:97:1--97:15. doi:10.4230/LIPIcs.ESA.2025.97.
\bibitem[Arya et al.(2004)Arya, Garg, Khandekar, Meyerson, Munagala, and Pandit]{Arya2004}
Arya V, Garg N, Khandekar R, Meyerson A, Munagala K, Pandit V (2004) Local search heuristics for $k$-median and facility location problems. \emph{SIAM Journal on Computing} 33(3):544--562. doi:10.1137/S0097539702416402.
\bibitem[Babaioff et al.(2022)Babaioff, Gonczarowski, and Nisan]{Babaioff2016}
Babaioff M, Gonczarowski YA, Nisan N (2022) The menu-size complexity of revenue approximation. \emph{Games and Economic Behavior} 134:281--307. doi:10.1016/j.geb.2021.03.001.
\bibitem[Beasley and Christofides(1989)]{BeasleyChristofides1989}
Beasley JE, Christofides N (1989) An algorithm for the resource constrained shortest path problem. \emph{Networks} 19(4):379--394. doi:10.1002/net.3230190402.
\bibitem[Bolusani et al.(2024)]{Bolusani2024}
Bolusani S, et al. (2024) The SCIP Optimization Suite 9.0. ZIB Report 24-02, Zuse Institute Berlin. The execution manifest records the version actually run; this citation documents the solver framework.
\bibitem[Di Puglia Pugliese and Guerriero(2013)]{PuglieseGuerriero2013}
Di Puglia Pugliese L, Guerriero F (2013) A survey of resource constrained shortest path problems: Exact solution approaches. \emph{Networks} 62(3):183--200. doi:10.1002/net.21511.
\bibitem[Downey and Fellows(1995)]{DowneyFellows1995b}
Downey RG, Fellows MR (1995) Fixed-parameter tractability and completeness II: On completeness for W[1]. \emph{Theoretical Computer Science} 141(1--2):109--131. doi:10.1016/0304-3975(94)00097-3.
\bibitem[Fisher(1981)]{Fisher1981}
Fisher ML (1981) The Lagrangian relaxation method for solving integer programming problems. \emph{Management Science} 27(1):1--18. doi:10.1287/mnsc.27.1.1.
\bibitem[Gurski et al.(2019)Gurski, Rehs, and Rethmann]{Gurski2019}
Gurski F, Rehs C, Rethmann J (2019) Knapsack problems: A parameterized point of view. \emph{Theoretical Computer Science} 775:93--108. doi:10.1016/j.tcs.2018.12.019.
\bibitem[Handler and Zang(1980)]{HandlerZang1980}
Handler GY, Zang I (1980) A dual algorithm for the constrained shortest path problem. \emph{Networks} 10(4):293--309. doi:10.1002/net.3230100403.
\bibitem[Hart and Nisan(2019)]{HartNisan2013}
Hart S, Nisan N (2019) Selling multiple correlated goods: Revenue maximization and menu-size complexity. \emph{Journal of Economic Theory} 183:991--1029. doi:10.1016/j.jet.2019.07.006.
\bibitem[Hassin(1992)]{Hassin1992}
Hassin R (1992) Approximation schemes for the restricted shortest path problem. \emph{Mathematics of Operations Research} 17(1):36--42. doi:10.1287/moor.17.1.36.
\bibitem[Karp(1972)]{Karp1972}
Karp RM (1972) Reducibility among combinatorial problems. Miller RE, Thatcher JW, eds. \emph{Complexity of Computer Computations} (Plenum Press, New York), 85--103. doi:10.1007/978-1-4684-2001-2\_9.
\bibitem[Lorenz and Raz(2001)]{LorenzRaz2001}
Lorenz DH, Raz D (2001) A simple efficient approximation scheme for the restricted shortest path problem. \emph{Operations Research Letters} 28(5):213--219. doi:10.1016/S0167-6377(01)00069-4.
\bibitem[Maher et al.(2016)Maher, Miltenberger, Pedroso, Rehfeldt, Schwarz, and Serrano]{Maher2016}
Maher S, Miltenberger M, Pedroso JP, Rehfeldt D, Schwarz R, Serrano F (2016) PySCIPOpt: Mathematical programming in Python with the SCIP optimization suite. \emph{Mathematical Software---ICMS 2016}, 301--307. doi:10.1007/978-3-319-42432-3\_37.
\bibitem[Mohajerin Esfahani and Kuhn(2018)]{EsfahaniKuhn2018}
Mohajerin Esfahani P, Kuhn D (2018) Data-driven distributionally robust optimization using the Wasserstein metric: Performance guarantees and tractable reformulations. \emph{Mathematical Programming} 171:115--166. doi:10.1007/s10107-017-1172-1.
\bibitem[Pag\`es and Wilbertz(2012)]{PagesWilbertz2012}
Pag\`es G, Wilbertz B (2012) Intrinsic stationarity for vector quantization: Foundation of dual quantization. \emph{SIAM Journal on Numerical Analysis} 50(2):747--780. doi:10.1137/110827041.
\bibitem[Patriksson(2008)]{Patriksson2008}
Patriksson M (2008) A survey on the continuous nonlinear resource allocation problem. \emph{European Journal of Operational Research} 185(1):1--46. doi:10.1016/j.ejor.2006.12.006.
\bibitem[Rusmevichientong et al.(2010)Rusmevichientong, Shen, and Shmoys]{Rusmevichientong2010}
Rusmevichientong P, Shen ZJM, Shmoys DB (2010) Dynamic assortment optimization with a multinomial logit choice model and capacity constraint. \emph{Operations Research} 58(6):1666--1680. doi:10.1287/opre.1100.0866.
\bibitem[Sadykov et al.(2021)Sadykov, Uchoa, and Pessoa]{Sadykov2021}
Sadykov R, Uchoa E, Pessoa A (2021) A bucket graph-based labeling algorithm with application to vehicle routing. \emph{Transportation Science} 55(1):4--28. doi:10.1287/trsc.2020.0985.
\bibitem[Virtanen et al.(2020)Virtanen, Gommers, Oliphant, et al.]{Virtanen2020}
Virtanen P, Gommers R, Oliphant TE, et al. (2020) SciPy 1.0: Fundamental algorithms for scientific computing in Python. \emph{Nature Methods} 17:261--272. doi:10.1038/s41592-019-0686-2.
\bibitem[Winkler(1988)]{Winkler1988}
Winkler G (1988) Extreme points of moment sets. \emph{Mathematics of Operations Research} 13(4):581--587. doi:10.1287/moor.13.4.581.
\bibitem[Wu(1991)]{Wu1991}
Wu X (1991) Optimal quantization by matrix searching. \emph{Journal of Algorithms} 12(4):663--673. doi:10.1016/0196-6774(91)90039-2.
\bibitem[Yang et al.(2025)Yang, Song, Wang, Yang, and Leus]{Yang2025}
Yang K, Song G, Wang R, Yang H, Leus R (2025) Perspective Benders decomposition with applications to fixed-charge nonlinear resource allocation. \emph{INFORMS Journal on Computing}, published online July 8. doi:10.1287/ijoc.2024.0984.
\bibitem[Zhang(2012)]{Zhang2012}
Zhang H (2012) Solving an infinite horizon adverse selection model through finite policy graphs. \emph{Operations Research} 60(4):850--864. doi:10.1287/opre.1120.1056.


\end{thebibliography}
\label{refs-end}\clearpage
\begin{table}[p]\centering
\caption{Antecedents, retained results, and the tariff-sensitive frontier}\label{tab:frontier60}
\begin{tabular}{p{0.24\textwidth}p{0.31\textwidth}p{0.36\textwidth}}\toprule
Problem or method & Relevant antecedent & Result and boundary in this paper\\\midrule
Resource-constrained paths & Lagrangian bounds, labeling and approximation; Handler and Zang (1980), Hassin (1992), Sadykov et al. (2021), Ahmadi et al. (2025) & Original-model reduction with endogenous arc allocations, exact capacity conservation and same-path equality repair. No generic new labeling method.\\[0.07in]
Fixed-charge nonlinear allocation & Perspective reformulation and Benders cuts; Yang et al. (2025) & Ordered terminal-capacity compression; exact polynomial optimization at uniform fees and fixed-parameter optimization in exceptional commands.\\[0.07in]
Parameterized packing & Knapsack parameters and kernel questions; Gurski et al. (2019) & W[1]-hardness in command allowance plus cap count for the original feasible-policy language; no membership or kernel claim.\\[0.07in]
Finite support and prototypes & Moment-set geometry, dual quantization and facility location; Winkler (1988), Pag\`es and Wilbertz (2012), Arya et al. (2004) & Retained cap-count bound, oscillation transfer and cap-preserving grouping frontier; generic support and clustering facts are antecedents.\\[0.07in]
Price certificates & Standard weak duality and fixed-book concavity & Same-book support characterizes equality of a price bound; distance bounds and a positive-gap uniform-fee example distinguish it from tariff tractability.\\\bottomrule
\end{tabular}
\end{table}
\begin{table}[p]\centering\caption{Controlled scaling: target attainment within the original total allowance}\label{tab:scaling63}
\begin{singlespace}\small\begin{tabular}{rrrrrrrrr}\toprule
$N$ & $k$ & $m$ & Sec. & Tariff & Price & Hybrid & Enumer. & SCIP\\\midrule
17 & 12 & 8 & 1 & 6/8 & 8/8 & 8/8 & 0/8 & 8/8 \\
17 & 12 & 8 & 6 & 8/8 & 8/8 & 8/8 & 0/8 & 8/8 \\
65 & 32 & 16 & 1 & 3/8 & 0/8 & 0/8 & 0/8 & 0/8 \\
65 & 32 & 16 & 6 & 5/8 & 8/8 & 8/8 & 0/8 & 8/8 \\
129 & 48 & 32 & 1 & 0/8 & 0/8 & 0/8 & 0/8 & 0/8 \\
129 & 48 & 32 & 6 & 3/8 & 0/8 & 3/8 & 0/8 & 0/8 \\
\bottomrule\end{tabular}\par\smallskip\begin{minipage}{\textwidth}\footnotesize Every cell includes all eight exception counts. Tariff, price, hybrid, and enumeration use independently checked rational upper and lower bounds. SCIP uses its numerical upper bound and an independently checked lower policy. A later successful replay does not change these original timing outcomes. All entries are constructed controlled cells, not independent random replicates.\end{minipage}\end{singlespace}\end{table}
\begin{table}[p]\centering\caption{The exception loop and independent certification at fixed physical size}\label{tab:d63}
\begin{singlespace}\small\begin{tabular}{rrrrrrr}\toprule
$d$ & Optimize (s) & Serialize (s) & Check (s) & Opt. KiB & Check KiB & Proof bytes\\\midrule
0 & 0.0051 & 0.0004 & 0.2088 & 27644 & 25716 & 4,143 \\
1 & 0.0061 & 0.0006 & 0.2110 & 27644 & 25776 & 5,674 \\
2 & 0.0070 & 0.0007 & 0.2102 & 27644 & 25720 & 8,699 \\
4 & 0.0120 & 0.0018 & 0.2163 & 27644 & 25736 & 25,368 \\
6 & 0.0266 & 0.0060 & 0.2359 & 27644 & 25720 & 84,248 \\
8 & 0.0782 & 0.0237 & 0.2964 & 27644 & 26628 & 308,314 \\
10 & 0.2457 & 0.0840 & 0.4932 & 33672 & 36360 & 1,119,607 \\
12 & 0.8189 & 0.2962 & 1.1157 & 55688 & 63496 & 3,964,865 \\
\bottomrule\end{tabular}\par\smallskip\begin{minipage}{\textwidth}\footnotesize Fixed $(N,k,m)=(17,12,8)$ and a six-second total allowance. Resident memory is the separate optimizer/checker process high-water mark. Proof bytes are uncompressed. Setup, process startup, and compression also consume the parent deadline; displayed component times are not asserted to sum to the total. The raw records contain all components and budgets.\end{minipage}\end{singlespace}\end{table}
\begin{table}[p]\centering\caption{Exact root-support classification and observed price-search work}\label{tab:root63}
\begin{singlespace}\small\begin{tabular}{lrrrrrr}\toprule
Root class & Cases & Med. splits & Max. splits & Max. depth & Med. solve (s) & Med. check (s)\\\midrule
Supported & 20 & 0 & 0 & 0 & 0.0028 & 0.0048 \\
Positive gap & 28 & 4 & 13 & 7 & 0.0314 & 0.0198 \\
\bottomrule\end{tabular}\par\smallskip\begin{minipage}{\textwidth}\footnotesize Classification comes from complete feasible-book enumeration and the exact upper concave hull, not from a selected probe book. All 48 final certificates are exact; 1,176 books are enumerated in total. Root-gap cases have a maximum exact minimum-price gap of $1/20$. Local diagnostic times are descriptive and are not pooled with the earlier frozen execution environment.\end{minipage}\end{singlespace}\end{table}
\end{document}
